\documentclass{article}
\usepackage[utf8]{inputenc}
\usepackage{graphicx}
\usepackage{url}
\usepackage[T1]{fontenc}
\usepackage{hyperref}
\usepackage{listings}
\usepackage{xcolor}
\usepackage{float}
\usepackage[inkscapelatex=false]{svg}
\usepackage{algorithm}
\usepackage{algpseudocode}
\usepackage{amssymb}

\lstset{
  basicstyle=\ttfamily\footnotesize, 
  numbers=left,                     
  numberstyle=\tiny\color{gray},
  keywordstyle=\color{blue},        
  commentstyle=\color{green!50!black},
  stringstyle=\color{orange},        
  frame=single,                     
  breaklines=true,                   
  tabsize=2,
  showstringspaces=false,
  columns=flexible,
  frame=single,
  backgroundcolor=\color{gray!5}
}

\hypersetup{
    colorlinks=true,
    linkcolor=black,
    urlcolor=blue,
    citecolor=blue
}

\begin{document}

\begin{titlepage}
    \centering
    \vspace*{2cm}

    % Título principal
    {\Huge \bfseries Segunda Entrega\par}
    \vspace{1.5cm}

    % Autores
    {\Large Gonzalo Pérez\par}
    \vspace{1cm}

    % Curso / asignatura
    {\large Ingeniería Informática\\
        Universidad del País Vasco (EHU) \\
        Desarrollo Avanzado de Software\par}
    \vspace{2cm}

    % Fecha
    {\large Abril 2026\par}

    \vfill
\end{titlepage}


\tableofcontents
\newpage


\section{Introducción}

En esta segunda entrega se sigue con la aplicación iniciada en la primera entrega. En el enunciado de esta entrega se han propuesto nuevos requisitos y funcionalidades a implementar.

\begin{itemize}
    \item Uso de bases de datos remotas.
    \item Integración con Google Maps y Geolocalización.
    \item Capturar imagenes desde la camara, guardarlas en el servidor y mostrarlas en la aplicación.
\end{itemize}

Adicionalmente, se ha proporcionado una lista de características opcionales para enriquecer la aplicación:

\begin{itemize}
    \item Uso de algún content provider.
    \item Servicio en primer plano y mensajes broadcast.
    \item Mensajería FCM.
    \item Widget con actualizacion periodica.
    \item Uso de servicio o tarea programada mediante alarma.
\end{itemize}

Para la implementación de estos requisitos se ha seguido la implementación de la aplicación de la primera entrega, añadiendo las nuevas funcionalidades requeridas.

\section{Funcionalidades implementadas}


\subsection{Elementos obligatorios}


\subsubsection{Base de datos remota}

Se ha implementado un servidor con una base de datos remota para la gestión de tareas y usuarios, proporcionando un sistema de persistencia más robusto que el almacenamiento local. Para mitigar la dependencia de la conectividad a internet, se ha desarrollado un mecanismo de sincronización entre ambos entornos. \\

Bajo esta arquitectura, las operaciones de escritura (creación, modificación o borrado) se ejecutan prioritariamente en la base de datos remota y se replican posteriormente en la local. Por el contrario, las consultas de lectura se realizan directamente sobre la base de datos local para minimizar la latencia de red y optimizar la respuesta de la interfaz. Adicionalmente, el sistema ejecuta una actualización automática cada 15 minutos, garantizando así la integridad y consistencia de los datos almacenados en el dispositivo.

\subsubsection{Google Maps y Geolocalización}

Se ha integrado la funcionalidad de asignar una ubicación geográfica a las tareas. En consecuencia, se ha incorporado un acceso directo en las tarjetas de cada tarea que permite abrir la aplicación Google Maps con las coordenadas preseleccionadas. Esta característica centraliza la información relevante en una única interfaz, optimizando la experiencia del usuario al eliminar la necesidad de buscar manualmente la dirección en aplicaciones externas.

\subsubsection{Captura de imágenes y gestión de perfiles}

Se ha integrado una funcionalidad que permite al usuario capturar una fotografía para su perfil personal. El proceso de gestión de archivos se divide en dos niveles: el archivo de imagen se transfiere al servidor para su almacenamiento en un directorio interno del sistema de archivos, mientras que en la base de datos se registra únicamente la URI correspondiente.

\subsection{Elementos opcionales}


\subsubsection{Content Provider}

Se ha integrado el \textit{Content Provider} de Google Calendar para permitir la sincronización de las tareas creadas con el calendario del usuario. Esta funcionalidad es opcional y puede gestionarse directamente desde el menú de ajustes de la aplicación.\\

Paralelamente, se ha desarrollado un \textit{Content Provider} propio que actúa como capa de abstracción para la base de datos local. Este componente centraliza todas las operaciones CRUD de la aplicación y permite que aplicaciones externas, con los permisos adecuados, interactúen con los datos de las tareas. Internamente, la aplicación utiliza este proveedor como intermediario exclusivo para cualquier consulta o modificación de los datos locales, garantizando la integridad de la información.

\subsubsection{Servicio en primer plano y mensajes \textit{broadcast}}

Se ha desarrollado un servicio en primer plano (\textit{Foreground Service}) encargado de monitorizar la ubicación del usuario en tiempo real, manteniendo una notificación persistente que informa al usuario sobre el estado de la monitorización. \\

Este servicio trabaja en conjunto con la funcionalidad de geolocalización para emitir alertas de proximidad. Mediante el uso de \textit{Broadcast Receivers}, el sistema detecta variaciones en las coordenadas geográficas y, al identificar que el usuario se encuentra en un radio cercano a una tarea pendiente, envía una señal para activar la notificación correspondiente.

\subsubsection{Mensajería mediante Firebase Cloud Messaging (FCM)}

Se ha implementado un sistema de notificaciones basado en \textit{Firebase Cloud Messaging} para informar a los usuarios sobre el lanzamiento de nuevas versiones de la aplicación. El flujo de trabajo se ha automatizado mediante un entorno de CI/CD utilizando GitHub Actions (\href{https://github.com/GonzaloPerezGomez/DASPROYECTO/blob/das-002/.github/workflows/notify_new_version.yml}{workflow}). \\

Al detectar un incremento en la versión de la aplicación (clave \texttt{versionCode} en el archivo \texttt{build.gradle}), se dispara un flujo de trabajo que envía automáticamente un mensaje FCM al canal o «topic» denominado \texttt{'nueva\_version'}. De este modo, todos los dispositivos suscritos reciben una notificación en tiempo real sin intervención manual del desarrollador.

\subsubsection{Widget}

Se ha implementado un widget que muestra las 3 proximas tareas pendientes del usuario. El widget se actualiza automaticamente cada vez que se añade o modifica una tarea, o cada 30 minutos. Esto permite al usuario tener la informacion mas relevante sobre sus tareas sin necesidad de abrir la aplicación.

\subsubsection{Servicio programado mediante alarma}

La aplicacion cuenta con un sistema de notificaciones que informa todos los dias a las 8 de la mañana sobre las tareas del dia o de dias anteriores que no se han completado. Este servicio se implemento en la primera entrega y no se ha modificado.

\newpage


\section{Clases y Base de datos}


\subsection{Arquitectura y Organización de Paquetes}

Siguiendo el enfoque modular de la anterior entrega se han reorganizado las clases por responsabilidades.

\begin{itemize}
    \item \textbf{ui}: Contiene todas las clases relacionadas con la interfaz de usuario, subdivididas en \texttt{activities}, \texttt{fragments}, \texttt{adapters} (para el RecyclerView), \texttt{dialogs} y \texttt{views} (preferencias personalizadas).
    \item \textbf{data}: Centraliza la gestión de datos, incluyendo la persistencia local con Room (\texttt{db}) y el \texttt{\textit{Content Provider}} propio (\texttt{provider}).
    \item \textbf{services}: Aloja los procesos en segundo plano, diferenciando entre servicios de sistema (\texttt{background}) como FCM o proximidad, y tareas asíncronas de WorkManager (\texttt{workers}).
    \item \textbf{receiver}: Agrupa los \texttt{\textit{Broadcast Receivers}} que capturan eventos del sistema, como el arranque del dispositivo o las alarmas programadas.
    \item \textbf{utils}: Clases auxiliares para la gestión del idioma y otras utilidades comunes.
    \item \textbf{widget}: Contiene la lógica específica para el funcionamiento y actualización del widget de escritorio.
\end{itemize}

Esto mejora mucho la legibilidad del código y facilita el mantenimiento del mismo. Esto se ha visto especialmente necesario en esta entrega debido a la cantidad de funcionalidades y clases nuevas que se han añadido.

\subsection{Base de Datos y Persistencia Local (Room)}

Para esta entrega se ha migrado el sistema de almacenamiento local a \textbf{Room Persistence Library}. Esta capa de abstracción sobre SQLite no solo mejora la robustez del código mediante la verificación de consultas en tiempo de compilación, sino que facilita la gestión de hilos y la reactividad de la interfaz.

La implementación se basa en tres pilares fundamentales:
\begin{itemize}
    \item \textbf{Entity (\texttt{TareaEntity})}: Define la estructura de la tabla \texttt{tareas}, incluyendo anotaciones para claves primarias y mapeo de columnas que coinciden con el esquema remoto.
    \item \textbf{DAO (\texttt{TareaDao})}: Interfaz que centraliza las consultas SQL. Se ha hecho uso de \texttt{SimpleSQLiteQuery} para permitir el ordenamiento dinámico de las tareas (por fecha o prioridad) directamente desde la base de datos.
    \item \textbf{Database (\texttt{AppDatabase})}: Clase que mantiene la instancia única de la base de datos siguiendo el patrón \textit{Singleton}.
\end{itemize}

\begin{figure}[H]
    \centering
    \includesvg[width=0.7\textwidth]{img/esquema_db.svg}
    \caption{Diagrama de la base de datos remota y local.}
    \label{fig:esquema_db}
\end{figure}

\subsection{Patrón Repository y Single Source of Truth (SSOT)}

Se ha implementado la clase \texttt{TareaRepository} siguiendo el principio de \textbf{Única Fuente de Verdad}. El repositorio es el encargado de decidir de dónde proceden los datos y cómo se sincronizan:
\begin{itemize}
    \item \textbf{Lecturas}: Se realizan exclusivamente desde Room. Esto garantiza que la aplicación sea extremadamente rápida y funcional incluso cuando el dispositivo no tiene conexión a internet.
    \item \textbf{Escrituras}: Cuando el usuario crea o modifica una tarea, el repositorio intenta primero realizar la operación en el servidor remoto (MariaDB). Si la operación tiene éxito, se dispara automáticamente una sincronización para actualizar la caché local de Room.
\end{itemize}

\subsection{Content Provider e Interoperabilidad}

Cumpliendo con los requisitos del hito, se ha desarrollado un \texttt{\textit{Content Provider}} propio (\texttt{TareasContentProvider}). Este componente expone los datos de la aplicación a través de \texttt{URIs} personalizadas.

Internamente, el \textit{provider} delega todas las peticiones al \texttt{TareaRepository}, asegurando que cualquier aplicación tercera que interactúe con nuestros datos respete la misma lógica de negocio y sincronización que la aplicación principal.

\subsection{Sincronización y Procesos en Segundo Plano}

La comunicación con el servidor remoto se gestiona de forma asíncrona para no bloquear el hilo de la interfaz de usuario:
\begin{itemize}
    \item \textbf{WorkManager}: Se utilizan \texttt{SyncWorker} y \texttt{ConexionWorker} para tareas periódicas y comprobaciones de conectividad.
    \item \textbf{Comunicación HTTP}: Se realizan peticiones \texttt{POST} codificadas en \texttt{URL-form-encoded} hacia una API en PHP, procesando las respuestas en formato \texttt{JSON} para su posterior mapeo a objetos de Room.
\end{itemize}

En la Fig\ref{fig:flujo_db} se detalla el flujo completo de sincronización entre el cliente y el servidor.

\begin{figure}[H]
    \centering
    \makebox[\textwidth][c]{
        \includesvg[width=1.35\textwidth]{img/flujo_db.svg}
    }
    \caption{Diagrama de flujo de todo el sistema de persistencia y sincronización.}
    \label{fig:flujo_db}
\end{figure}
\subsection{Servicios en Segundo Plano y Notificaciones}

Para garantizar una experiencia proactiva, la aplicación hace uso de diversos servicios que operan independientemente de la visibilidad de la interfaz:
\begin{itemize}
    \item \textbf{Servicio de Proximidad (\texttt{ProximidadService})}: Implementado como un \textit{Foreground Service}, mantiene una ejecución persistente mediante una notificación en la barra de estado. Utiliza la API \texttt{FusedLocationProviderClient} de Google para monitorizar la ubicación del usuario. Cuando detecta que el dispositivo se encuentra en un radio de 200 metros de una tarea pendiente, emite un \textit{broadcast} interno y lanza una notificación de alta prioridad.
    \item \textbf{Mensajería Cloud (\texttt{MiFirebaseMessagingService})}: Integrado con \textit{Firebase Cloud Messaging} (FCM), este servicio está a la escucha de mensajes enviados desde el servidor. Está configurado para suscribirse al tema \textit{'nueva\_version'}, permitiendo que la aplicación reaccione a actualizaciones automáticas enviadas desde el flujo de CI/CD de GitHub Actions. Adicionalmente se puede realizar una demo del servicio accediendo a \url{http://34.68.1.253:81/fcm_enviar.php}. Es un formulario que permite definir el titulo, el cuerpo y el destino de la notificación.
\end{itemize}

\subsection{Widget de Escritorio}

Se ha desarrollado un widget (\texttt{TareasWidgetProvider}) que permite al usuario visualizar sus tareas más urgentes directamente desde la pantalla de inicio del dispositivo. Técnicamente, este componente se basa en las siguientes claves:
\begin{itemize}
    \item \textbf{Carga asíncrona}: Dado que Room no permite consultas en el hilo principal, el widget utiliza un \texttt{ExecutorService} para recuperar de forma segura las 3 próximas tareas pendientes.
    \item \textbf{Interfaz Remota (\texttt{RemoteViews})}: Debido a que los widgets se ejecutan en un proceso separado (el \textit{launcher}), la interfaz se gestiona mediante \texttt{RemoteViews}.
    \item \textbf{Actualización Reactiva}: El repositorio de datos se encarga de notificar al \texttt{AppWidgetManager} cada vez que ocurre una sincronización o un cambio local, asegurando que el widget siempre refleje la información más reciente.
\end{itemize}

\subsection{Diagrama de clases}

Para facilitar la comprensión de la estructura del proyecto, se han elaborado tres diagramas de clases UML, agrupados por capa funcional.

\begin{figure}[H]
    \centering
    \includesvg[width=0.9\textwidth]{img/diagrama1_actividades}
    \caption{Diagrama de clases: Actividades y navegación.}
\end{figure}

\begin{figure}[H]
    \centering
    \includesvg[width=0.9\textwidth]{img/diagrama2_fragmentos}
    \caption{Diagrama de clases: Fragmentos, adaptadores e interfaces.}
\end{figure}

\begin{figure}[H]
    \centering
    \includesvg[width=\textwidth]{img/Diagrama3_Datos_Dialogos}
    \caption{Diagrama de clases: Datos, diálogos y notificaciones.}
\end{figure}

\newpage

\section{Manual de Usuario}

\subsection{Login y Registro}

Al entrar en la aplicación, se muestra una pantalla de login. Si no se tiene un cuenta se puede ir a la pantalla de registro pulsando el botón "Registrarse". En caso de tener una cuenta se puede iniciar sesión introduciendo el correo electrónico y la contraseña.

\begin{figure}[H]
    \begin{minipage}{0.45\textwidth}
        \centering
        \includegraphics[width=0.8\textwidth]{img/login.jpeg}
        \caption{Pantalla de login.}
    \end{minipage}\hfill
    \begin{minipage}{0.45\textwidth}
        \centering
        \includegraphics[width=0.8\textwidth]{img/registro.jpeg}
        \caption{Pantalla de registro.}
    \end{minipage}
\end{figure}

Adicionalmente se puede marcar la casilla de "Recuerdame" para que en caso de cerrar sesión se mantenga el correo electrónico introducido. Se puede cerrar sesión desde el menu de ajustes.

\begin{figure}[H]
    \begin{minipage}{0.45\textwidth}
        \centering
        \includegraphics[width=0.8\textwidth]{img/cerrar_sesion.jpeg}
        \caption{Ajustes: Cerrar sesión.}
    \end{minipage}
\end{figure}

\subsection{Ubicación}

Como se ha comentado, se puede añadir una ubicación a una tarea. Esto se realiza pulsando el botón de "Añadir ubicacion" en el formulario de añadir o editar tarea.

\begin{figure}[H]
    \begin{minipage}{0.45\textwidth}
        \centering
        \includegraphics[width=0.6\textwidth]{img/seleccionar_ubicacion.jpeg}
        \caption{Seleccionar ubicación.}
    \end{minipage}
    \begin{minipage}{0.45\textwidth}
        \centering
        \includegraphics[width=0.6\textwidth]{img/new_tarea_ubi.jpeg}
        \caption{Añadir ubicación.}
    \end{minipage}
\end{figure}

Aquellas tareas que tengan una ubicación asignada tendrán un icono de un mapa en la lista de tareas, y pulsandolo se abrirá un mapa con la ubicación de la tarea.

\begin{figure}[H]
    \begin{minipage}{0.45\textwidth}
        \centering
        \includegraphics[width=0.6\textwidth]{img/con_sin_ubi.jpeg}
        \caption{Lista de tareas con ubicación.}
    \end{minipage}
    \begin{minipage}{0.45\textwidth}
        \centering
        \includegraphics[width=0.6\textwidth]{img/google_maps.jpeg}
        \caption{Mapa con la ubicación de la tarea.}
    \end{minipage}
\end{figure}

\subsection{Servicio de proximidad}

Desde el Navigation Drawer se puede iniciar el sistema de proximidad, el cual mostrará una notificación cuando el usuario se encuentre a 200 metros de una tarea con ubicación.

\begin{figure}[H]
    \begin{minipage}{0.45\textwidth}
        \centering
        \includegraphics[width=0.6\textwidth]{img/iniciar_proximidad.jpeg}
        \caption{Servicio de proximidad.}
    \end{minipage}
    \begin{minipage}{0.45\textwidth}
        \centering
        \includegraphics[width=0.6\textwidth]{img/not_prox_perma.jpeg}
        \caption{Notificación de proximidad.}
    \end{minipage}
\end{figure}

\subsection{Mensajería FCM}

Dentro de la aplicacion la mensajeria FCM se usa para notificar a los usuarios sobre una nueva version de la aplicacion. Estas notificaciones son completamente autonomas. Adicionalmente se pueden enviar notificaciones de forma manual, tanto a todos los usuarios como a usuarios individuales (\href{http://34.68.1.253:81/fcm_enviar.php}{demo}).

\begin{figure}[H]
    \begin{minipage}{0.45\textwidth}
        \centering
        \includegraphics[width=1.2\textwidth]{img/form_g.png}
        \caption{Formulario para enviar notificaciones.}
    \end{minipage}
    \begin{minipage}{0.45\textwidth}
        \centering
        \includegraphics[width=0.6\textwidth]{img/not_g.jpeg}
        \caption{Notificación FCM.}
    \end{minipage}
\end{figure}

\begin{figure}[H]
    \begin{minipage}{0.45\textwidth}
        \centering
        \includegraphics[width=1.2\textwidth]{img/form_todos.png}
        \caption{Formulario para enviar notificaciones.}
    \end{minipage}
    \begin{minipage}{0.45\textwidth}
        \centering
        \includegraphics[width=0.6\textwidth]{img/not_todos_usuarios.jpeg}
        \caption{Notificación FCM.}
    \end{minipage}
\end{figure}

\begin{figure}[H]
    \begin{minipage}{0.45\textwidth}
        \centering
        \includegraphics[width=0.6\textwidth]{img/not_nueva_version.jpeg}
        \caption{Notificación FCM.}
    \end{minipage}
\end{figure}

\subsection{Widget}

Un simplre widget que muestra las tareas pendientes.

\begin{figure}[H]
    \begin{minipage}{0.45\textwidth}
        \centering
        \includegraphics[width=0.6\textwidth]{img/widget.jpeg}
        \caption{Widget de tareas pendientes.}
    \end{minipage}
\end{figure}

\newpage

\section{Dificultades}

El desafío más significativo del proyecto radicó en la migración integral del sistema de persistencia hacia su arquitectura actual. La mayor complejidad consistió en definir una estructura organizada donde cada componente tuviera una responsabilidad clara, garantizando que la interacción entre los diferentes entornos de datos operase con total fluidez.

\section{Fuentes utilizadas}

Para la realización de este proyecto se han consultado diversas fuentes de documentación y herramientas de apoyo, entre las que destacan:

\begin{itemize}
    \item \textbf{Diapositivas de la asignatura: } Principal fuente de información para la realización del proyecto.
    \item \textbf{Documentación oficial de Android Developers:} Consultas varias puntuales(\url{https://developer.android.com/}).
    \item \textbf{Stack Overflow:} Utilizado puntualmente para la resolución de errores específicos y búsqueda de enfoques alternativos de implementación.(\url{https://stackoverflow.com/})
    \item \textbf{Asistente de IA (Antigravity):} Herramienta de inteligencia artificial usada para correción de errores, depuración de código, uso de buenas practicas de programación y revisión de la memoria.(\url{https://antigravity.google/})
\end{itemize}

\section{Repositorio de código}

El código fuente del proyecto se encuentra disponible en el siguiente repositorio de GitHub:

\begin{center}
    \url{https://github.com/GonzaloPerezGomez/DASPROYECTO/tree/das-002}
\end{center}

En la carpeta \texttt{server} se encuentran todos los ficheros php que se usan en el servidor remoto.

\end{document}